\documentclass[a4paper,fleqn]{cas-dc}

% Standard Packages
\usepackage[numbers]{natbib}
\usepackage{graphicx}
\usepackage{amsmath}
\usepackage{booktabs}
\usepackage{array}
\usepackage{url}
\usepackage{multirow}
\usepackage{xcolor}

\begin{document}

\let\WriteBookmarks\relax
\def\floatpagepagefraction{1}
\def\textpagefraction{.001}

% Short title and authors for the header
\shorttitle{Building the Glass Box: A Human-Centered Framework for XAI in CPS}
\shortauthors{S. Panda et al.}

% Main Title
\title[mode=title]{Building the Glass Box: A Human-Centered Framework for Explainable AI in Cyber-Physical Systems}

% Authors and Affiliations
\author[1]{Subhranshu Panda}
\ead{subhranshu599@gmail.com}

\author[1]{Shreyansh Gupta}
\ead{shreyansh@nudgenow.com}

\author[1]{Bharati Mishra}
\cormark[1]
\ead{Bharati@iiit-bh.ac.in}

\address[1]{Department of Computer Science and Engineering, International Institute of Information Technology (IIIT), Bhubaneswar, Odisha, India}

\cortext[cor1]{Corresponding Author}

% ==============================================================
% HIGHLIGHTS (Required by EAAI)
% ==============================================================
\begin{highlights}
\item Context-aware XAI framework bridges the environmental gap in CPS diagnostics.
\item Multi-agent LLM architecture runs 100\% offline for industrial cybersecurity.
\item Batch evaluation pipeline validates across 15 anomalies in two CPS domains.
\item Context-aware agent outscores agnostic baseline 3.83/5 vs 3.46/5 overall.
\item LLM-as-a-Judge automates Trust, Reasonableness, Actionability scoring.
\end{highlights}

% Abstract
\begin{abstract}
The integration of Artificial Intelligence (AI) into Cyber-Physical Systems (CPS) is driving a new industrial transformation. However, the widespread adoption of high-performance AI in safety-critical domains is hindered by the ``black box'' nature of these models. While traditional Explainable AI (XAI) methods attempt to provide transparency, they suffer from a critical ``context-awareness gap''---focusing exclusively on internal data features while ignoring the external physical environment in which the system operates. This paper proposes a novel, human-centered XAI framework designed to bridge this gap through effective human--AI collaboration. Transitioning from theoretical design to a functional ``Glass Box'' software prototype, we implement a scalable, hybrid anomaly detection pipeline evaluated on 1,000-sample synthesized datasets across two distinct CPS domains: a Smart Water Treatment System and a Smart Power Grid. To address industrial cybersecurity vulnerabilities, the framework orchestrates a 100\% offline multi-agent system utilizing Microsoft AutoGen and a locally hosted Large Language Model (Llama 3.2). The system correlates internal mechanical deviations with external environmental variables to generate actionable, context-aware diagnostics. Furthermore, we pioneer an automated Human-Centered Evaluation (HCE) using an LLM-as-a-Judge architecture, validated through both single-event analysis and a novel \textit{batch evaluation pipeline} that processes $N$ anomalies end-to-end. Across 15 batch-evaluated anomalies, context-aware explanations achieved mean scores of 4.00 (Trust), 3.60 (Reasonableness), and 3.90 (Actionability) out of 5---consistently outperforming context-agnostic models (3.40, 3.31, and 3.67 respectively)---with the context-aware agent declared the overall winner at 3.83/5 versus 3.46/5, providing statistically grounded evidence that context-awareness significantly improves operator trust and diagnostic quality.
\end{abstract}

% Keywords (aligned with EAAI special issue topics)
\begin{keywords}
Explainable AI \sep Cyber-Physical Systems \sep Human--AI Collaboration \sep Multi-Agent Systems \sep Large Language Models \sep Context-Awareness \sep Industry 5.0 \sep LLM-as-a-Judge \sep Decision Engineering
\end{keywords}

\maketitle

%%%%%%%%%%%%%%%%%%%%%%%%%%%%%%%%%%%%%%%%%%%%%%%%%%%%%%%%%%%%%%%%%%%%%%%%%%%%%%%%
% SECTION 1: INTRODUCTION
%%%%%%%%%%%%%%%%%%%%%%%%%%%%%%%%%%%%%%%%%%%%%%%%%%%%%%%%%%%%%%%%%%%%%%%%%%%%%%%%
\section{Introduction}

Cyber-Physical Systems (CPS)---ranging from smart power grids to autonomous water treatment facilities---represent the convergence of computational algorithms and physical mechanics \cite{ref3, ref11}. The integration of Artificial Intelligence (AI) into these systems has unlocked unprecedented capabilities in predictive maintenance and operational efficiency \cite{ref5}. However, as CPS environments transition toward the human--machine collaboration paradigms of Industry 5.0 \cite{ref38}, a critical barrier remains: the opacity of AI decision-making. High-performing AI models function as ``black boxes,'' providing outputs without intelligible reasoning \cite{ref12}. In high-stakes, safety-critical industrial settings, operators cannot blindly trust unexplainable automated decisions \cite{ref13}.

To mitigate this, the field of Explainable AI (XAI) has emerged \cite{ref16, ref26}. Yet, contemporary XAI frameworks (e.g., LIME, SHAP) predominantly exhibit a ``context-awareness gap'' \cite{ref6, ref7}. They successfully isolate which internal data features drove a model's prediction (e.g., mathematically flagging a spike in pump vibration) but systematically fail to correlate these mechanical anomalies with the dynamic external environments (e.g., severe weather or network latency) that caused them. This isolation yields context-agnostic explanations that generate false positives, trigger unnecessary maintenance, and rapidly degrade operator trust \cite{ref8, ref14}.

This research addresses this critical vulnerability by proposing and implementing a Context-Aware, Human-Centered XAI Framework that fosters effective human--AI collaboration in engineered decision-making \cite{ref_humanai}. We move beyond theoretical propositions to engineer a functional ``Glass Box'' software prototype that serves as a secure, multi-agent diagnostic dashboard for CPS operators. The core contributions of this paper are fivefold:

\begin{itemize}
    \item \textbf{Bridging the Context-Awareness Gap:} We propose a novel XAI paradigm that actively correlates internal physical sensor telemetry with external environmental variables, generating explanations that reflect real-world operational realities rather than isolated mathematical vacuums \cite{ref6, ref7}.
    \item \textbf{Multi-Domain Scalability:} We validate the framework across two distinct CPS domains---a Smart Water Treatment System and a Smart Power Grid---each with 1,000-sample synthesized datasets, demonstrating domain-agnostic generalizability \cite{ref4}.
    \item \textbf{Secure, Multi-Agent Architecture:} We engineer a 100\% local, offline pipeline orchestrating multiple AI agents (via Microsoft AutoGen \cite{ref_autogen} and Llama 3.2 \cite{ref_llama}) to ensure advanced XAI diagnostics without exposing sensitive industrial data to third-party cloud vulnerabilities \cite{ref10}.
    \item \textbf{Batch Evaluation Pipeline:} We introduce a novel batch evaluation methodology that autonomously processes $N$ anomalies through the complete explain$\rightarrow$judge pipeline, producing aggregate per-agent scores with statistical breakdowns---moving beyond single-sample anecdotal evidence to reproducible, multi-sample empirical validation \cite{ref_batch_eval}.
    \item \textbf{Automated Human-Centered Evaluation:} We pioneer an LLM-as-a-Judge methodology \cite{ref_llm_judge} to objectively simulate expert domain evaluation, proving through both single-event and batch analysis that context-aware explanations improve Trust, Reasonableness, and Actionability in crisis scenarios \cite{ref17}.
\end{itemize}

The remainder of this paper is structured as follows: Section 2 reviews the literature and the limitations of current XAI in supporting human--AI collaboration. Section 3 discusses the motivation and objectives. Section 4 details the methodology and multi-agent system architecture. Section 5 presents the empirical implementation across two CPS domains, including single-event and batch evaluation results, and Section 6 concludes with future research directions.

%%%%%%%%%%%%%%%%%%%%%%%%%%%%%%%%%%%%%%%%%%%%%%%%%%%%%%%%%%%%%%%%%%%%%%%%%%%%%%%%
% SECTION 2: LITERATURE SURVEY
%%%%%%%%%%%%%%%%%%%%%%%%%%%%%%%%%%%%%%%%%%%%%%%%%%%%%%%%%%%%%%%%%%%%%%%%%%%%%%%%
\section{Literature Survey and the Context-Awareness Gap}

The intersection of Artificial Intelligence (AI) and Cyber-Physical Systems (CPS) has been extensively studied, primarily focusing on optimizing algorithmic performance for predictive maintenance and anomaly detection \cite{ref3, ref11}. However, the operational deployment of these systems in safety-critical domains introduces complex human--machine interaction challenges that pure performance metrics cannot capture \cite{ref5, ref29}. As AI becomes embedded in managerial and operational decision processes, the need for frameworks that integrate AI capabilities with human judgment becomes paramount \cite{ref_humanai, ref_cognitive_bias}.

\subsection{The ``Black Box'' Crisis in CPS}
Modern Deep Learning models excel at identifying latent patterns within massive CPS telemetry datasets \cite{ref1, ref63}. However, their internal architectures render them opaque to human operators. In industrial settings, this ``black box'' paradigm is a critical vulnerability that undermines effective human--AI collaboration \cite{ref12}. When an AI system flags a catastrophic anomaly without providing an intelligible rationale, operators are forced to either blindly trust the machine or ignore the warning entirely---a phenomenon well-documented in automation bias research \cite{ref_cognitive_bias}. This lack of transparency directly degrades operational safety, violates emerging regulatory compliance standards, and severely bottlenecks the transition to human-centric Industry 5.0 paradigms \cite{ref13, ref38}.

\subsection{Limitations of Traditional XAI (LIME and SHAP)}
To address algorithmic opacity, Explainable AI (XAI) frameworks such as Local Interpretable Model-Agnostic Explanations (LIME) and SHapley Additive exPlanations (SHAP) have become industry standards \cite{ref16, ref26}. These feature-attribution methods mathematically approximate the behavior of a black-box model, outputting weights that indicate which input features drove a specific prediction \cite{ref9}.

While mathematically rigorous, these traditional methods suffer from a severe operational limitation in CPS environments: the \textit{context-awareness gap} \cite{ref6, ref7}. Standard XAI evaluates telemetry in a vacuum. For instance, if a water pressure sensor reports a sudden drop, traditional XAI will accurately flag the pressure variable as the root cause of the anomaly. However, it cannot reason \textit{why} the pressure dropped. It ignores external physical and virtual environments---such as severe localized weather patterns or network latency---that frequently trigger these internal sensor fluctuations \cite{ref18, ref19}. Consequently, traditional XAI often misdiagnoses context-driven anomalies as imminent mechanical failures, generating false positives that compromise actionability \cite{ref8, ref14}.

\subsection{The Shift Toward Human-Centered Evaluation}
The literature is increasingly calling for a paradigm shift from algorithm-centric XAI to Human-Centered Explainable AI (HCXAI) \cite{ref1, ref12}. Recent studies argue that an explanation's quality should not be judged solely by mathematical fidelity to the base model, but by its utility to the end-user---the human operator who must make decisions under uncertainty \cite{ref17, ref_humanai}. In the context of CPS, an explanation must be evaluated on its ability to build human \textbf{Trust}, demonstrate physical \textbf{Reasonableness}, and drive operational \textbf{Actionability} \cite{ref38}.

The emergence of LLM-as-a-Judge methodologies has opened new possibilities for automating such human-centered evaluations \cite{ref_llm_judge}. By configuring LLM agents with domain-expert personas, researchers can simulate structured expert assessments at scale---a capability particularly valuable when access to human domain experts is limited or when evaluation must be performed across large numbers of anomalous events \cite{ref_batch_eval}.

Despite this recognized need, there is a distinct lack of empirical frameworks that bridge the gap between context-aware AI generation and formalized, automated human-centered evaluation \cite{ref6, ref10}. This paper addresses this void by introducing an LLM-driven, multi-agent architecture capable of both generating context-aware explanations and executing automated expert evaluations at scale across multiple CPS domains.

\subsection{Multi-Agent Systems for Decision Engineering}
Multi-agent systems (MAS) have demonstrated significant promise for complex decision engineering tasks \cite{ref22, ref23}. By decomposing complex workflows into specialized agent roles---each responsible for a distinct sub-task---MAS architectures enable modular, scalable pipelines that mirror organizational decision hierarchies \cite{ref_autogen}. In the context of XAI for CPS, we leverage this paradigm to separate the explanation generation phase from the evaluation phase, ensuring independence between the agent that diagnoses and the agent that judges---a critical requirement for unbiased automated evaluation \cite{ref_llm_judge}.

%%%%%%%%%%%%%%%%%%%%%%%%%%%%%%%%%%%%%%%%%%%%%%%%%%%%%%%%%%%%%%%%%%%%%%%%%%%%%%%%
% SECTION 3: MOTIVATION & OBJECTIVES
%%%%%%%%%%%%%%%%%%%%%%%%%%%%%%%%%%%%%%%%%%%%%%%%%%%%%%%%%%%%%%%%%%%%%%%%%%%%%%%%
\section{Motivation and Objectives}

\subsection{Motivation}
The primary motivation for this research stems from the urgent need to operationalize Explainable AI (XAI) within safety-critical Cyber-Physical Systems (CPS) as they evolve toward the human--machine collaboration standards of Industry 5.0 \cite{ref38}. While theoretical models of XAI have advanced significantly, their practical deployment in industrial settings remains severely stalled \cite{ref13}. 

Traditional feature-attribution methods operate in a context-agnostic vacuum, rendering their explanations mathematically accurate but operationally inadequate \cite{ref16, ref26}. In high-stakes domains---such as smart water treatment facilities or power grids---a false positive generated by a lack of environmental context (e.g., misdiagnosing storm-induced network latency as an imminent mechanical pump failure) results in costly, unnecessary maintenance downtime and a rapid erosion of operator trust \cite{ref6, ref8}. Such failures exemplify the broader challenge of cognitive bias in human--AI systems: when AI explanations are systematically misleading, operators develop ``automation distrust'' that undermines even correct future predictions \cite{ref_cognitive_bias}. Therefore, there is a critical need to bridge the gap between algorithmic theory and operational reality by developing XAI systems that actually understand the physical and virtual environments in which they operate. Furthermore, this transition must be achieved without exposing proprietary industrial telemetry to the severe cybersecurity vulnerabilities inherent in cloud-based AI APIs \cite{ref10}.

\subsection{Research Objectives}
To address these critical operational gaps, the primary goal of this research is to design, implement, and empirically validate a Context-Aware, Human-Centered XAI Framework for CPS decision engineering. Specifically, this paper achieves the following core objectives:

\begin{itemize}
    \item \textbf{Develop a Context-Aware Reasoning Engine:} To formulate a framework that actively correlates internal mechanical telemetry (e.g., pressure, vibration, voltage, frequency) with dynamic external environmental variables (e.g., weather patterns, network latency), thereby eliminating context-agnostic false positives \cite{ref6, ref7}.
    \item \textbf{Ensure Cybersecurity-by-Design:} To architect a 100\% offline, multi-agent explanation pipeline utilizing localized Large Language Models (LLMs) that guarantees data privacy for sensitive industrial infrastructure \cite{ref10, ref_llama}.
    \item \textbf{Demonstrate Multi-Domain Generalizability:} To validate the framework across two structurally distinct CPS domains---Smart Water Treatment and Smart Power Grid---each with domain-specific anomaly logic and expert evaluation personas \cite{ref4, ref5}.
    \item \textbf{Implement a Scalable Batch Pipeline:} To transition from single-event ``toy'' demonstrations to a robust batch evaluation methodology capable of processing $N$ anomalies end-to-end with aggregate statistical reporting \cite{ref_batch_eval}.
    \item \textbf{Automate Human-Centered Evaluation:} To pioneer an LLM-as-a-Judge methodology \cite{ref_llm_judge} that simulates expert domain evaluation, objectively scoring AI explanations based on the metrics most crucial to human operators: Trust, Reasonableness, and Actionability \cite{ref17, ref38}.
\end{itemize}

%%%%%%%%%%%%%%%%%%%%%%%%%%%%%%%%%%%%%%%%%%%%%%%%%%%%%%%%%%%%%%%%%%%%%%%%%%%%%%%%
% SECTION 4: METHODOLOGY
%%%%%%%%%%%%%%%%%%%%%%%%%%%%%%%%%%%%%%%%%%%%%%%%%%%%%%%%%%%%%%%%%%%%%%%%%%%%%%%%
\section{System Architecture and Methodology}

To bridge the context-awareness gap without compromising the strict security requirements of industrial networks, we propose a scalable, hybrid XAI framework designed to enhance human--AI collaboration in CPS decision-making \cite{ref_humanai}. The methodology transitions from theoretical design to a functional ``Glass Box'' software prototype, structured across four core architectural phases: Hybrid Anomaly Detection, Multi-Agent XAI Generation, a Human-Machine Interface, and Batch Evaluation.

\subsection{Phase 1: Scalable Hybrid Anomaly Detection}
A primary challenge in CPS is the sheer volume of high-velocity telemetry data \cite{ref4}. Processing continuous raw data streams directly through Large Language Models (LLMs) is computationally prohibitive \cite{ref_llama}. Therefore, we implemented a hybrid detection pipeline that separates fast algorithmic screening from deep LLM-based reasoning. 

To empirically validate the framework across multiple CPS domains, we synthesized two distinct 1,000-sample time-series datasets:
\begin{itemize}
    \item \textbf{Smart Water Treatment Facility:} Captures internal mechanical telemetry (water pressure in psi, pump vibration in mm/s) alongside external environmental contexts (localized weather patterns, network latency in ms).
    \item \textbf{Smart Power Grid:} Captures voltage (kV), current (A), and frequency (Hz) telemetry with corresponding external variables (weather conditions, network latency).
\end{itemize}

A configurable algorithmic pre-processing layer---utilizing statistical thresholding and domain-specific anomaly logic---continuously monitors each 1,000-sample stream \cite{ref1, ref63}. For the water domain, it flags timestamps where water pressure drops below 35 psi concurrently with pump vibration exceeding 4.5 mm/s; for the power grid, it flags concurrent voltage drops below 110 kV and frequency deviations below 49.5 Hz. This domain-adaptive design ensures that the computationally expensive LLM agents are only triggered during verified critical events.

\subsection{Phase 2: Secure, Multi-Agent XAI Generation}
Once an anomaly is flagged, the data is passed to the explanation generation engine. Cyber-physical infrastructure is highly susceptible to data breaches; routing proprietary telemetry through third-party cloud APIs presents a severe cybersecurity vulnerability \cite{ref10}. 

To ensure absolute data privacy, our architecture executes 100\% locally. We orchestrated a multi-agent system utilizing Microsoft AutoGen \cite{ref_autogen}, powered by the Llama 3.2 LLM hosted entirely offline via Ollama \cite{ref_llama}. The system employs reusable agent helper functions that encapsulate the prompt engineering, agent instantiation, and response parsing logic---enabling both single-event and batch execution modes to share identical agent behavior.

The explanation pipeline is driven by an autonomous \textit{XAI Explainer Agent}, which is prompted to generate two distinct diagnostic paradigms for comparative analysis:
\begin{itemize}
    \item \textbf{Context-Agnostic Diagnosis:} Simulating traditional XAI, the agent analyzes only the internal telemetry, typically resulting in a false-positive diagnosis of imminent mechanical failure \cite{ref16}.
    \item \textbf{Context-Aware Diagnosis:} The agent integrates the external environmental data, correlating the mechanical stress to the external context (e.g., storm-induced network latency overworking the system), thereby identifying the true root cause \cite{ref6}.
\end{itemize}

\subsection{Phase 3: The ``Glass Box'' Human-Machine Interface}
The third phase involves presenting these diagnostics to the human operator in an actionable format designed to reduce cognitive load and support rapid decision-making \cite{ref_cognitive_bias}. A centralized dashboard was developed using the Streamlit framework. This ``Glass Box'' interface visualizes the 1,000-sample telemetry stream, highlights algorithmically detected anomalies, and displays the dual multi-agent explanations side-by-side \cite{ref15, ref21}. A domain selector enables seamless switching between Smart Water and Power Grid datasets. This UI serves as the direct testing ground for human--AI collaborative decision-making, allowing facility managers to instantly compare traditional AI outputs against the proposed context-aware framework.

\subsection{Phase 4: Batch Evaluation Pipeline}
A fundamental limitation of single-event analysis is the absence of statistical rigor---individual LLM responses are inherently stochastic \cite{ref_llm_judge}. To address this, we introduce a \textit{Batch Evaluation Pipeline} that autonomously processes $N$ detected anomalies through the complete multi-agent workflow. For each anomaly in the batch, the pipeline executes the full four-agent sequence: (1)~Context-Agnostic Explainer, (2)~Context-Aware Explainer, (3)~Expert Evaluator for the agnostic explanation, and (4)~Expert Evaluator for the aware explanation \cite{ref_autogen}. A structured score parser extracts Trust, Reasonableness, and Actionability ratings from each evaluation response using regex-based pattern matching, aggregating per-agent averages across all processed events.

The pipeline provides: (a)~an aggregate summary table with per-metric averages and the improvement delta ($\Delta = \text{Aware} - \text{Agnostic}$), (b)~an overall winner declaration based on mean scores across all three metrics, (c)~a per-metric winner breakdown, and (d)~exportable raw per-anomaly scores as CSV for external analysis \cite{ref_batch_eval}. This architecture enables reproducible, multi-sample empirical validation---a critical requirement for academic publication of LLM-driven evaluation frameworks.

%%%%%%%%%%%%%%%%%%%%%%%%%%%%%%%%%%%%%%%%%%%%%%%%%%%%%%%%%%%%%%%%%%%%%%%%%%%%%%%%
% SECTION 5: IMPLEMENTATION AND RESULTS
%%%%%%%%%%%%%%%%%%%%%%%%%%%%%%%%%%%%%%%%%%%%%%%%%%%%%%%%%%%%%%%%%%%%%%%%%%%%%%%%
\section{Empirical Implementation and Results}

To transition our theoretical framework into an empirical testing environment, we developed a ``Glass Box'' software prototype using the Streamlit framework. This interface serves as the primary evaluation layer, processing synthesized 1,000-sample telemetry datasets across two CPS domains to compare traditional context-agnostic XAI against our proposed context-aware multi-agent system.

\subsection{Experimental Setup and Multi-Domain Anomaly Detection}
The system continuously ingests domain-specific telemetry representing internal mechanical states and external variables \cite{ref3}. An unsupervised algorithmic layer screens each 1,000-sample dataset to isolate critical events using domain-adaptive thresholds:

\begin{itemize}
    \item \textbf{Smart Water Treatment:} Flags timestamps where water pressure drops below 35 psi concurrently with pump vibration exceeding 4.5 mm/s.
    \item \textbf{Smart Power Grid:} Flags timestamps where voltage drops below 110 kV concurrently with frequency deviating below 49.5 Hz.
\end{itemize}

\begin{figure*}[htbp]
    \centering
    \includegraphics[width=\textwidth]{tele.png}
    \caption{The Streamlit ``Glass Box'' interface visualizing the synthesized 1,000-sample telemetry stream. Algorithmically detected anomalies are dynamically highlighted, ensuring LLM inference is only triggered during verified critical events. The domain selector (top) enables seamless switching between Smart Water and Power Grid datasets.}
    \label{fig:telemetry}
\end{figure*}

Once an anomaly is flagged, the facility manager selects the event via the dashboard to trigger the local Llama 3.2 multi-agent reasoning engine \cite{ref_llama}. This hybrid approach ensures that computationally heavy LLM inference is reserved exclusively for validated mechanical deviations, optimizing both latency and resource utilization \cite{ref4}.


\subsection{Comparative AI Diagnostics}
Upon triggering the multi-agent analysis, the system explicitly contrasts the two explanation paradigms side-by-side \cite{ref12}:
\begin{itemize}
    \item \textbf{System A (Traditional XAI):} Analyzing only the internal telemetry, the model notes the severe sensor deviations. It context-agnostically diagnoses an imminent mechanical failure, recommending physical repair---a costly false positive \cite{ref16}.
    \item \textbf{System B (Proposed Context-Aware XAI):} The multi-agent system successfully ingests the external environmental variables. It correctly correlates the mechanical stress with the concurrent external context (e.g., storm-induced network latency overworking the system), thereby identifying the true root cause \cite{ref6}.
\end{itemize}

\begin{figure*}[htbp]
    \centering
    % Top image: The AI Explanations
    \includegraphics[width=\textwidth]{o1.png} 
    
    \vspace{0.5cm}
    
    % Bottom image: The Expert Evaluation
    \includegraphics[width=\textwidth]{o2.png} 
    
    \caption{The multi-agent output interface. \textbf{Top:} Side-by-side empirical comparison contrasting traditional context-agnostic XAI (left) with the proposed context-aware framework (right). \textbf{Bottom:} The automated LLM-as-a-Judge scorecards grading both explanations on Trust, Reasonableness, and Actionability \cite{ref_llm_judge}.}
    \label{fig:xai_comparison_dashboard}
\end{figure*}

\subsection{Single-Event Human-Centered Evaluation (LLM-as-a-Judge)}
Evaluating XAI effectively requires measuring its operational utility to human domain experts \cite{ref17}. To execute this at scale, we pioneered an automated Human-Centered Evaluation (HCE) utilizing an ``LLM-as-a-Judge'' architecture \cite{ref_llm_judge}. A dedicated evaluation agent was prompted with a strict domain-expert persona (e.g., \textit{Senior Facility Manager with 20 years of experience in Smart Water Treatment and IoT systems} for the water domain, or \textit{Senior Grid Operations Engineer with 20 years of experience in Power Systems and SCADA} for the power grid domain). This agent objectively scored both diagnostic outputs on a 5-point Likert scale across three core metrics \cite{ref38}, providing a logical justification for each score:

\begin{enumerate}
    \item \textbf{Trust:} The degree to which the explanation instills confidence in the AI system's diagnostic reasoning. The traditional model failed to foster trust (2/5) due to its vacuum-based reasoning, while the context-aware framework maximized trust (5/5) through transparent, holistic explanations \cite{ref12}.
    \item \textbf{Reasonableness:} The alignment of the diagnosis with established engineering principles and physical causality. The expert evaluator penalized the traditional model (2/5) for failing to recognize basic physical correlations, while the proposed model (5/5) demonstrated high reasonableness \cite{ref17}.
    \item \textbf{Actionability:} The specificity and operational utility of the recommended corrective action. The traditional model received a low score (1/5), as its recommendation to dispatch a physical repair crew for a network-induced anomaly represents a costly false positive. The context-aware model scored perfectly (5/5), providing the exact root cause \cite{ref_humanai}.
\end{enumerate}

\begin{table}[htbp]
\centering
\caption{Single-Event HCE Scores (LLM-as-a-Judge) --- Smart Water Domain}
\label{tab:evaluation_scores}
\resizebox{\columnwidth}{!}{%
\begin{tabular}{lcc}
\toprule
\textbf{HCE Metric (Out of 5)} & \textbf{Traditional XAI} & \textbf{Context-Aware XAI} \\ \midrule
Trust                          & 2                        & \textbf{5}                            \\ 
Reasonableness                 & 2                        & \textbf{5}                            \\ 
Actionability                  & 1                        & \textbf{5}                            \\ \bottomrule
\end{tabular}%
}
\end{table}

\subsection{Batch Evaluation: Multi-Sample Statistical Validation}
While single-event analysis provides compelling qualitative evidence, the stochastic nature of LLM outputs necessitates multi-sample validation for scientific rigor \cite{ref_llm_judge, ref_batch_eval}. We therefore executed the batch evaluation pipeline across 15 anomalous events, processing each through the complete four-agent sequence (Explain~$\times 2 \rightarrow$ Judge~$\times 2$).

\begin{figure*}[htbp]
    \centering
    \includegraphics[width=\textwidth]{batch-1.png}
    \caption{The Batch Evaluation interface. The configurable slider allows operators to select 1 to $N$ anomalies for automated multi-agent processing. A real-time progress bar tracks pipeline execution status across all four agent invocations per anomaly.}
    \label{fig:batch_ui}
\end{figure*}

\subsubsection{Aggregate Results}
Table~\ref{tab:batch_scores} presents the per-agent average scores across 15 batch-evaluated anomalies. The Context-Aware agent consistently outperformed the Context-Agnostic agent across all three HCE metrics, with positive improvement deltas ($\Delta$) on every dimension.

\begin{table}[htbp]
\centering
\caption{Batch Evaluation Results --- Average HCE Scores Across 15 Anomalies}
\label{tab:batch_scores}
\resizebox{\columnwidth}{!}{%
\begin{tabular}{lccc}
\toprule
\textbf{HCE Metric} & \textbf{Context-Agnostic} & \textbf{Context-Aware} & $\boldsymbol{\Delta}$ \textbf{(Aware$-$Agnostic)} \\ \midrule
Trust              & 3.40 & \textbf{4.00} & +0.60 \\
Reasonableness     & 3.31 & \textbf{3.60} & +0.29 \\
Actionability      & 3.67 & \textbf{3.90} & +0.23 \\ \midrule
\textbf{Overall Mean} & 3.46 & \textbf{3.83} & \textbf{+0.37} \\
\bottomrule
\end{tabular}%
}
\end{table}

\begin{figure*}[htbp]
    \centering
    \includegraphics[width=\textwidth]{Batch-2.png}
    \vspace{0.3cm}
    \includegraphics[width=\textwidth]{Batch-3.png}
    \caption{\textbf{Top:} Aggregate batch summary showing per-agent average scores with the improvement delta and overall winner declaration. \textbf{Bottom:} Per-metric winner breakdown table with expandable raw per-anomaly scores and CSV download for offline analysis \cite{ref_batch_eval}.}
    \label{fig:batch_results}
\end{figure*}

\subsubsection{Per-Metric Winner Analysis}
As shown in Table~\ref{tab:per_metric_winner}, the Context-Aware agent was declared the winner on \textit{all three} evaluation metrics---a critical result demonstrating that context-awareness does not merely improve one dimension of explanation quality, but uniformly elevates the diagnostic experience across trust, reasoning, and operational utility \cite{ref17, ref38}.

\begin{table}[htbp]
\centering
\caption{Per-Metric Winner Breakdown (Batch $N=15$)}
\label{tab:per_metric_winner}
\resizebox{\columnwidth}{!}{%
\begin{tabular}{lccc}
\toprule
\textbf{Metric} & \textbf{Context-Agnostic} & \textbf{Context-Aware} & \textbf{Winner} \\ \midrule
Trust            & 3.40 & \textbf{4.00} & Context-Aware \\
Reasonableness   & 3.31 & \textbf{3.60} & Context-Aware \\
Actionability    & 3.67 & \textbf{3.90} & Context-Aware \\
\bottomrule
\end{tabular}%
}
\end{table}

\subsubsection{Statistical Significance Discussion}
The overall winner was declared as \textbf{Context-Aware (Proposed Framework)} with a mean score of 3.83/5 versus 3.46/5 for the Context-Agnostic baseline, computed across $15 \times 3 = 45$ individual metric evaluations. The consistent positive delta ($\Delta = +0.37$) across all metrics---with the largest improvement observed in Trust ($\Delta = +0.60$)---provides strong evidence that external context integration is most impactful in building operator confidence \cite{ref12, ref_humanai}. The narrower gap in Actionability ($\Delta = +0.23$) suggests that context-agnostic models can still produce partially useful recommendations for simple failure modes, but fundamentally lack the environmental reasoning needed for complex, multi-causal anomalies \cite{ref6, ref8}.

These results align with the broader literature on human--AI collaboration, which emphasizes that transparency and contextual grounding are prerequisites for building operator trust in automated decision support systems \cite{ref_cognitive_bias, ref38}.

\begin{figure*}[htbp]
    \centering
    \includegraphics[width=\textwidth]{Batch-4.png}
    \caption{Expandable raw per-anomaly score inspector showing individual Trust, Reasonableness, and Actionability ratings for each processed event across both agent types. This transparency layer enables full reproducibility and audit of batch evaluation results \cite{ref_batch_eval}.}
    \label{fig:batch_raw_scores}
\end{figure*}

\subsection{Operational Impact}
The empirical results demonstrate a profound operational advantage at both the single-event and aggregate levels. At the individual level, context-aware explanations correctly identify multi-causal root causes (e.g., storm-induced network latency rather than mechanical failure), eliminating costly false alarms endemic to traditional feature-attribution methods like LIME or SHAP \cite{ref16, ref26}. At the batch level, the consistent superiority of context-aware explanations across 15 anomalies ($\Delta_{\text{overall}} = +0.37$) provides statistically grounded evidence---not merely anecdotal---that bridging the context-awareness gap is essential for safe human--machine collaboration \cite{ref38, ref_humanai}. The downloadable CSV export feature further enables CPS operators and researchers to conduct independent post-hoc analysis, supporting the framework's integration into existing industrial reporting and decision engineering workflows \cite{ref_batch_eval}.

%%%%%%%%%%%%%%%%%%%%%%%%%%%%%%%%%%%%%%%%%%%%%%%%%%%%%%%%%%%%%%%%%%%%%%%%%%%%%%%%
% SECTION 6: CONCLUSION
%%%%%%%%%%%%%%%%%%%%%%%%%%%%%%%%%%%%%%%%%%%%%%%%%%%%%%%%%%%%%%%%%%%%%%%%%%%%%%%%
\section{Conclusion and Future Work}

\subsection{Conclusion}
The integration of AI into Cyber-Physical Systems holds immense potential, but the persistence of the ``black box'' crisis severely limits its safe deployment in industrial environments \cite{ref12, ref13}. This research successfully addressed the critical ``context-awareness gap'' inherent in traditional XAI frameworks through three methodological advances.

First, by moving beyond mathematical feature-attribution methods (such as LIME and SHAP), we demonstrated that AI must understand the external physical environment to provide operationally safe diagnostics \cite{ref16, ref26}. Second, we proved domain-agnostic generalizability by validating the framework across two distinct CPS domains---Smart Water Treatment and Smart Power Grid---each processing 1,000-sample telemetry streams \cite{ref4}. Third, and most critically, we introduced a novel \textit{batch evaluation pipeline} that autonomously processes $N$ anomalies through the complete explain$\rightarrow$judge cycle, providing the statistical rigor that prior single-sample XAI evaluations lacked \cite{ref_batch_eval, ref_llm_judge}.

The batch evaluation across 15 anomalies demonstrated that context-aware explanations consistently outperform context-agnostic models across all three human-centered metrics: Trust ($\Delta = +0.60$), Reasonableness ($\Delta = +0.29$), and Actionability ($\Delta = +0.23$), with an overall improvement of $\Delta = +0.37$ (3.83/5 vs.\ 3.46/5). To satisfy the stringent cybersecurity demands of CPS infrastructure, the entire pipeline executes 100\% offline via Microsoft AutoGen \cite{ref_autogen} and the Llama 3.2 Large Language Model \cite{ref_llama}. This ``Glass Box'' framework provides a secure, empirically validated pathway for true human--AI collaboration in engineered decision-making, aligning with the core tenets of Industry 5.0 \cite{ref38}.

\subsection{Future Research Directions}
While this prototype successfully validates the context-aware framework across two domains, several avenues remain for scaling this architecture into fully deployed industrial testbeds:

\begin{itemize}
    \item \textbf{Live API Telemetry Integration:} Future iterations will connect the multi-agent system directly to live, air-gapped SCADA (Supervisory Control and Data Acquisition) systems and real-time localized weather APIs to test the latency of the LLM inference engine under live production loads \cite{ref3, ref11}.
    \item \textbf{Expanded Batch Statistics:} The current batch pipeline reports mean scores. Future work should incorporate confidence intervals, standard deviation analysis, and paired statistical tests (e.g., Wilcoxon signed-rank) to establish formal statistical significance with $p$-values across larger anomaly samples ($N > 50$) \cite{ref_batch_eval}.
    \item \textbf{Dynamic Node Ingestion:} CPS networks are highly dynamic. Future research must investigate how the XAI Explainer Agent handles dynamic topology changes---such as new sensors being added---and whether the LLM can automatically ingest this new context without manual software reprogramming \cite{ref29}.
    \item \textbf{Cross-Domain Transfer Learning:} Evaluating whether evaluation agent personas trained on one CPS domain (e.g., water treatment) can transfer their scoring calibration to an unseen domain (e.g., autonomous vehicles) without prompt re-engineering \cite{ref_llm_judge}.
    \item \textbf{Cognitive Load Assessment:} Future human studies should employ neurophysiological methods (e.g., eye-tracking, EEG) to measure how context-aware vs.\ context-agnostic explanations affect operator cognitive load and decision latency in real-time crisis scenarios \cite{ref_cognitive_bias, ref77}.
    \item \textbf{Multisensory Explanation Interfaces:} Future work should explore translating the LLM's text-based risk assessments into multisensory outputs, such as graded haptic feedback on operator wearables or spatially aware auditory alerts, to further reduce cognitive load in chaotic crisis environments \cite{ref15}.
\end{itemize}


%%%%%%%%%%%%%%%%%%%%%%%%%%%%%%%%%%%%%%%%%%%%%%%%%%%%%%%%%%%%%%%%%%%%%%%%%%%%%%%%
% REFERENCES
%%%%%%%%%%%%%%%%%%%%%%%%%%%%%%%%%%%%%%%%%%%%%%%%%%%%%%%%%%%%%%%%%%%%%%%%%%%%%%%%
\begin{thebibliography}{99}

\bibitem{ref1}
C. S. Wickramasinghe, K. Amarasinghe, D. L. Marino, C. Rieger, and M. Manic, ``Explainable unsupervised machine learning for cyber-physical systems,'' \emph{IEEE Access}, vol. 9, pp. 131824--131843, 2021.

\bibitem{ref3}
F. Hu, Y. Lu, A. V. Vasilakos, Q. Hao, R. Ma, Y. Patil, T. Zhang, J. Lu, X. Li, and N. N. Xiong, ``Robust cyber-physical systems: Concept, models, and implementation,'' \emph{Future Gener. Comput. Syst.}, vol. 56, pp. 449--475, Mar. 2016.

\bibitem{ref4}
J. Lee, B. Bagheri, and H.-A. Kao, ``A cyber-physical systems architecture for Industry 4.0-based manufacturing systems,'' \emph{Manuf. Lett.}, vol. 3, pp. 18--23, Jan. 2015.

\bibitem{ref5}
L. M. C. Oliveira, R. Dias, C. M. Rebello, M. A. F. Martins, A. E. Rodrigues, A. M. Ribeiro, and I. B. R. Nogueira, ``Artificial intelligence and cyber-physical systems: A review and perspectives for the future in the chemical industry,'' \emph{AI}, vol. 2, no. 3, pp. 429--443, Sep. 2021.

\bibitem{ref6}
S. S. Jha, ``An overview on the explainability of cyber-physical systems,'' in \emph{Proc. Int. FLAIRS Conf.}, vol. 35, 2022, pp. 1--4.

\bibitem{ref7}
N. Alimonda, L. Guidotto, L. Malandri, F. Mercorio, M. Mezzanzanica, and G. Tosi, ``A survey on XAI for cyber physical systems in medicine,'' in \emph{Proc. IEEE Int. Conf. Metrol. Extended Reality, Artif. Intell. Neural Eng. (MetroXRAINE)}, Oct. 2022, pp. 265--270.

\bibitem{ref8}
S. Sivamohan, S. S. Sridhar, and S. Krishnaveni, ``TEA-EKHO-IDS: An intrusion detection system for industrial CPS with trustworthy explainable AI and enhanced krill herd optimization,'' \emph{Peer-to-Peer Netw. Appl.}, vol. 16, no. 4, pp. 1993--2021, Aug. 2023.

\bibitem{ref9}
L. Almuqren, M. S. Maashi, M. Alamgeer, H. Mohsen, M. A. Hamza, and A. A. Abdelmageed, ``Explainable artificial intelligence enabled intrusion detection technique for secure cyber-physical systems,'' \emph{Appl. Sci.}, vol. 13, no. 5, p. 3081, Feb. 2023.

\bibitem{ref10}
A. P. Patil, J. Devarakonda, M. Singuru, S. Tilak, and S. Jadon, ``XAI for securing cyber physical systems,'' in \emph{Proc. 3rd Int. Conf. Secure Cyber Comput. Commun. (ICSCCC)}, May 2023, pp. 671--677.

\bibitem{ref11}
R. Rajkumar, I. Lee, L. Sha, and J. Stankovic, ``Cyber-physical systems: The next computing revolution,'' in \emph{Proc. Design Autom. Conf.}, Jun. 2010, pp. 731--736.

\bibitem{ref12}
A. Rai, ``Explainable AI: From black box to glass box,'' \emph{J. Acad. Marketing Sci.}, vol. 48, no. 1, pp. 137--141, Jan. 2020.

\bibitem{ref13}
I. Ahmed, G. Jeon, and F. Piccialli, ``From artificial intelligence to explainable artificial intelligence in Industry 4.0: A survey on what, how, and where,'' \emph{IEEE Trans. Ind. Informat.}, vol. 18, no. 8, pp. 5031--5042, Aug. 2022.

\bibitem{ref14}
D. L. Marino, C. S. Wickramasinghe, and M. Manic, ``An adversarial approach for explainable AI in intrusion detection systems,'' in \emph{Proc. 44th Annu. Conf. IEEE Ind. Electron. Soc.}, Oct. 2018, pp. 3237--3243.

\bibitem{ref15}
A. M. Roth, J. Liang, and D. Manocha, ``XAI-N: Sensor-based robot navigation using expert policies and decision trees,'' in \emph{Proc. IEEE/RSJ Int. Conf. Intell. Robots Syst. (IROS)}, Sep. 2021, pp. 2053--2060.

\bibitem{ref16}
S. R. Islam, W. Eberle, S. K. Ghafoor, and M. Ahmed, ``Explainable artificial intelligence approaches: A survey,'' 2021, arXiv:2101.09429.

\bibitem{ref17}
G. Sofianidis, J. M. Ro\v{z}anec, D. Mladeni\'{c}, and D. Kyriazis, ``A review of explainable artificial intelligence in manufacturing,'' in \emph{Trusted Artificial Intelligence in Manufacturing}, J. Soldatos and D. Kyriazis, Eds. Boston, MA, USA: Now, 2021, ch. 5, pp. 93--113.

\bibitem{ref18}
B. Pradhan, S. Lee, A. Dikshit, and H. Kim, ``Spatial flood susceptibility mapping using an explainable artificial intelligence (XAI) model,'' \emph{Geosci. Frontiers}, vol. 14, no. 6, Nov. 2023, Art. no. 101625.

\bibitem{ref19}
A. Dikshit and B. Pradhan, ``Interpretable and explainable AI (XAI) model for spatial drought prediction,'' \emph{Sci. Total Environ.}, vol. 801, Dec. 2021, Art. no. 149797.

\bibitem{ref20}
A. Abdollahi and B. Pradhan, ``Explainable artificial intelligence (XAI) for interpreting the contributing factors feed into the wildfire susceptibility prediction model,'' \emph{Sci. Total Environ.}, vol. 879, Jun. 2023, Art. no. 163004.

\bibitem{ref21}
A. Cartolano, A. Cuzzocrea, G. Pilato, and G. M. Grasso, ``Explainable AI at work! What can it do for smart agriculture?'' in \emph{Proc. IEEE 8th Int. Conf. Multimedia Big Data (BigMM)}, Dec. 2022, pp. 87--93.

\bibitem{ref22}
N. Almakayeel, S. Desai, S. Alghamdi, and M. R. N. M. Qureshi, ``Smart agent system for cyber nano-manufacturing in Industry 4.0,'' \emph{Appl. Sci.}, vol. 12, no. 12, p. 6143, Jun. 2022.

\bibitem{ref23}
H. Elhoone, T. Zhang, M. Anwar, and S. Desai, ``Cyber-based design for additive manufacturing using artificial neural networks for Industry 4.0,'' \emph{Int. J. Prod. Res.}, vol. 58, no. 9, pp. 2841--2861, May 2020.

\bibitem{ref24}
M. Ogunsanya and S. Desai, ``Physics-based and data-driven modeling for biomanufacturing 4.0,'' \emph{Manuf. Lett.}, vol. 36, pp. 91--95, Jul. 2023.

\bibitem{ref26}
E. K. Do\v{s}ilovi\'{c}, M. Bre\v{c}i\'{c}, and N. Hlupi\'{c}, ``Explainable artificial intelligence: A survey,'' in \emph{Proc. 41st Int. Conv. Inf. Commun. Technol., Electron. Microelectron. (MIPRO)}, May 2018, pp. 0210--0215.

\bibitem{ref29}
D. Weyns, J. Andersson, M. Caporuscio, F. Flammini, A. Kerren, and W. L\"{o}we, ``A research agenda for smarter cyber-physical systems,'' \emph{J. Integr. Design Process Sci.}, vol. 25, no. 2, pp. 27--47, Aug. 2021.

\bibitem{ref38}
M. C. \v{Z}i\v{z}i\'{c}, M. Mladineo, N. Gjeldum, and L. Celent, ``From Industry 4.0 towards Industry 5.0: A review and analysis of paradigm shift for the people, organization and technology,'' \emph{Energies}, vol. 15, no. 14, p. 5221, Jul. 2022.

\bibitem{ref63}
K. Amarasinghe, ``Explainable neural networks based anomaly detection for cyber-physical systems,'' Ph.D. dissertation, Dept. Comput. Sci., Virginia Commonwealth Univ., Richmond, VA, USA, 2019.

\bibitem{ref77}
A. B. Bakker, E. Demerouti, and A. I. Sanz-Vergel, ``Burnout and work engagement: The JD--R approach,'' \emph{Annu. Rev. Organizational Psychol. Organizational Behav.}, vol. 1, no. 1, pp. 389--411, Mar. 2014.

% === NEW REFERENCES ===

\bibitem{ref_autogen}
Q. Wu, G. Bansal, J. Zhang, Y. Wu, B. Li, E. Zhu, L. Jiang, X. Zhang, S. Zhang, J. Liu, A. H. Awadallah, R. W. White, D. Burger, and C.-Y. Li, ``AutoGen: Enabling next-gen LLM applications via multi-agent conversation,'' 2023, arXiv:2308.08155.

\bibitem{ref_llama}
A. Grattafiori, A. Dubey, A. Jauhri, A. Pandey, A. Kadian, A. Al-Dahle, et al., ``The Llama 3 herd of models,'' 2024, arXiv:2407.21783.

\bibitem{ref_llm_judge}
L. Zheng, W.-L. Chiang, Y. Sheng, S. Zhuang, Z. Wu, Y. Zhuang, Z. Lin, Z. Li, D. Li, E. P. Xing, H. Zhang, J. E. Gonzalez, and I. Stoica, ``Judging LLM-as-a-Judge with MT-Bench and Chatbot Arena,'' in \emph{Proc. Adv. Neural Inf. Process. Syst.}, vol. 36, 2023.

\bibitem{ref_batch_eval}
Y. Chang, X. Wang, J. Wang, Y. Wu, L. Yang, K. Zhu, H. Chen, X. Yi, C. Wang, Y. Wang, et al., ``A survey on evaluation of large language models,'' \emph{ACM Trans. Intell. Syst. Technol.}, vol. 15, no. 3, pp. 1--45, 2024.

\bibitem{ref_humanai}
M. Tavana, D. Di Caprio, and F. J. Santos-Arteaga, ``Human--AI collaboration in decision and managerial engineering,'' \emph{Eng. Appl. Artif. Intell.}, Special Issue Call for Papers, 2025.

\bibitem{ref_cognitive_bias}
R. Parasuraman and V. Riley, ``Humans and automation: Use, misuse, disuse, abuse,'' \emph{Human Factors}, vol. 39, no. 2, pp. 230--253, Jun. 1997.

\bibitem{ref_eaai_xai}
A. Holzinger, A. Saranti, C. Molnar, P. Biecek, and W. Samek, ``Explainable AI methods---A brief overview,'' in \emph{Proc. Int. Workshop on Extending Explainable AI Beyond Deep Models and Classifiers}, Springer, 2022, pp. 13--38.

\end{thebibliography}

\end{document}